\xmpheader 9/{Miscellany}
\TeX\ understand the following units for measuring distances:
\halign{\hskip 2\parindent \tt#\quad\hfil&#\hfil\cr
pt & point (72.27 points $=$ 1 inch)\cr
pc & pica (1 pica $=$ 12 points)\cr
in & inch (1 inch $=$ 2.54 centimeters)\cr
cm & centimeter (1 centimeter = 10 millimeters)\cr
mm & millimeter\cr
}
\medskip
By default, the left and top margins will be 1 inch, but this can be changed (see the source above).
\hoffset -5mm \voffset 2mm % move to the left by 5 mm and move down 2mm.
The ink goes in the rectangle $hsize \times vsize$.  The default value is $6 \times 8.9$ in inches.
\hsize = 6in \vsize=8.9in
\medskip
You can include figure in connection with the {\tt dvips}.  For that you need to include the 
macro package {\tt epsf.tex}.
\input epsf.tex
Now you can say
\vskip 12 pt plus 4 pt minus 4 pt % equivalent to \bigskip
\leftline{\hskip\parindent\epsfysize=2in\epsffile{picture.ps}}\bigskip
\noindent Picture is in the form of encapsulated PostScript.
